% HD-S50. GPT-6 Astra Pro /root/sixfolds; inherited effort, 21 November 2026.
\section{HD-S50: four-dimensional blocks for the rank-four normal subgroup}
\label{sec:four-block-kernel-scope}

Let \(K\subset\mathbf C\) have characteristic zero and let
\(N\subset GL_{4r}(K)\), \(r\ge6\), be finite and
absolutely irreducible, generated by actual elements of moving
rank four. Suppose that it preserves a complex system of
\(r\) blocks of dimension four. Then its kernel \(T\) on
these blocks is characteristic in \(N\), and every original
block projector belongs to \(M_{4r}(K)\). Every actual
\(K\)-linear normalizer of \(N\) preserves these same
\(K\)-blocks.

For the natural source from a connected full-monodromy sphere
cover, let
\(N\triangleleft G\subset GL_{4r}(K)\), with the exact
character field
\[
 K=\mathbf Q(i)\quad\hbox{or}\quad
 K=\mathbf Q(\sqrt{-d}),\qquad
 d\in\{2,7,11,19,43,67,163\}.                           \tag{S50.1}
\]
For the primitive Jacobian corner, restricted theta polarization,
signature \((3,3)\) and period rank nine, this actual block
system reduces the source to a faithful rank-four sphere
coefficient system with the same ordinary determinant class.
HD-S40's complete rank-at-most-fourteen source result then
applies: the Gaussian class is in \(\{[-1],[-3]\}\),
and the seven other fields have class \([-1]\) or no such
source. No assertion that every quotient inertia is a reflection
is used. The proof of characteristicness below uses no finite
\(GL_4\) classification.

\paragraph{1. The literal pure permutation subgroup and local generators.}
Call the original complex blocks \(V_1,\ldots,V_r\).
For a cycle of length \(\ell\), compose its actual
transition maps around the cycle, obtaining \(A:V_i\to V_i\).
Its fixed vectors are exactly the successive transports of
\(\ker(A-I)\), so its moving rank is
\[
             4\ell-\dim\ker(A-I)\ge4(\ell-1).
                                                               \tag{S50.2}
\]
A rank-four generator that permutes blocks must therefore
interchange exactly two blocks, with inverse transition maps,
and be the identity on every other block. In particular it is
an actual involution. The graph of these transposed pairs is
connected: its components would otherwise give invariant sums
of blocks for all the generators and hence for \(N\).

Choose a spanning tree, fix a basis in one block, and transport
it recursively along the chosen inverse transition maps on
its edges. In the resulting actual complex coordinate map,
the selected generators are the literal transpositions of
\((\mathbf C^4)^r\). They generate a pure copy of \(S_r\)
inside \(N\). For any other outside generator, multiply it
by the pure transposition of the same pair. The resulting
block-diagonal matrix is supported on that pair. A generator
fixing every block is supported on at most four blocks,
because every nonidentity block contributes moving rank at
least one. Let \(T_0\) be generated by all pure \(S_r\)
conjugates of these block-diagonal matrices. It is stable
under the pure permutations, and \(\langle T_0,S_r\rangle\)
contains every original generator. Thus, literally in these
coordinates,
\[
          N=T\rtimes S_r,\qquad T=T_0\subset GL_4(\mathbf C)^r.
                                                               \tag{S50.3}
\]
This proves both the split subgroup and the local generation
statement; no splitting of an arbitrary abstract extension is
being assumed.

The pure permutations identify all coordinate images of \(T\)
with the same finite subgroup \(L\subset GL_4(\mathbf C)\).
It is irreducible. Otherwise the direct sum of the same proper
\(L\)-invariant subspace in each coordinate would be a
proper nonzero subspace invariant under both factors in
(S50.3), hence under \(N\).

For a set \(A\) of four block labels write
\[
 T_A=\{t\in T:\ t_j=I\text{ for }j\notin A\}.
                                                               \tag{S50.4}
\]
For every \(i\in A\), the coordinate map \(T_A\to L\)
is onto. Indeed the entries at \(i\) of the support-at-most-four
generating matrices in (S50.3) generate \(L\). A pure
permutation fixing \(i\) moves the other labels in the
support of each such generator into \(A\), without changing
its entry at \(i\). All these entries therefore occur in
the subgroup image of \(T_A\).

\paragraph{2. Actual localization, including six blocks.}
If \(r\ge7\), choose two four-sets \(A,B\) with
\(A\cap B=\{i\}\). For \(a\in T_A\), \(b\in T_B\),
the actual block commutator is the identity outside \(i\)
and has entry \([a_i,b_i]\) at \(i\). The two coordinate
maps are onto \(L\), so their commutators generate the
actual localized copy of \(L'\). Therefore
\[
                         (L')^r\subset T\qquad(r\ge7).
                                                               \tag{S50.5}
\]

For \(r=6\), choose \(A,B\) of size four with
\(A\cap B=\{i,j\}\). Their commutators have support at
most two, and their entries at \(i\) include every
\([a_i,b_i]\), with \(a_i,b_i\in L\) arbitrary.
Let \(T_1\) be generated by all such first commutators,
for all such pairs \(A,B\), and all their conjugates by
\(T\) and the pure \(S_6\). Every one of these generators
has support at most two and all its entries lie in \(L'\).
Conjugation preserves those two facts. Its coordinate image
is exactly \(L'\): the preceding arbitrary commutators
generate \(L'\), while containment follows from the
normality of \(L'\) in \(L\).

More precisely, for each pair \(\{i,j\}\), the subgroup
of \(T_1\) supported on that pair maps onto \(L'\) at
\(i\). The original first commutators with intersection
that pair already supply generators for this image. Choose
\(j\ne k\), both different from \(i\). Commuting the
subgroups supported on \(\{i,j\}\) and \(\{i,k\}\)
now gives actual matrices supported only at \(i\), whose
entries generate \([L',L']=L''\). Hence
\[
                          (L'')^6\subset T\qquad(r=6).
                                                               \tag{S50.6}
\]
The use of subgroup images in both stages is essential:
products remain supported on the stated set. No arbitrary
choice of lifts is claimed to preserve support.

\paragraph{3. A universal characteristic-kernel argument.}
Put \(Q=L'\) for \(r\ge7\) and \(Q=L''\) for
\(r=6\), and write \(U=Q_1\cdots Q_r\) for the actual
coordinate copies furnished by (S50.5)--(S50.6). Then
\[
 U=Q^r\triangleleft N,\qquad U\subset T,
 \qquad T/U\lhook\joinrel\longrightarrow (L/Q)^r.
                                                               \tag{S50.7}
\]
The injection is induced by the coordinate quotients, whose
kernel on \(T\) is exactly \(Q^r\). Thus \(T/U\)
is solvable: \(L/L'\) is abelian, and \(L/L''\) has
abelian derived subgroup.

Consider any surjective homomorphism
\(\varphi:N\to S_r\). The subgroup
\(R=\varphi(U)\) is normal in \(S_r\), and is therefore
\(1,A_r\), or \(S_r\), since \(r\ge5\).
For clarity, this normal-subgroup fact follows from simplicity
of \(A_r\): a normal subgroup meeting it trivially
centralizes it, and its centralizer in \(S_r\) is trivial.
The latter assertion follows from the point stabilizers, whose
only fixed labels are the respective points.

The subgroups \(R_i=\varphi(Q_i)\) commute pairwise,
are normal in \(R\), generate \(R\), and are permuted
transitively by conjugation from \(\varphi(N)=S_r\).
If \(R\ne1\), its normal subgroups imply that one
\(R_i\) contains \(A_r\): for \(R=A_r\) this is
simplicity, and for \(R=S_r\) it is the normal-subgroup
list just proved. Transitivity gives the same containment
for every \(R_i\). Two of them cannot commute, since the
centralizer of \(A_r\) in \(S_r\) is trivial. Hence
\(\varphi(U)=1\). It follows that \(\varphi(T)\)
is a solvable normal subgroup of \(S_r\), so is trivial.

For any automorphism \(\beta\) of \(N\), apply this
result to the surjection \(\pi\beta:N\to S_r\), where
\(\pi\) is the actual block action. Then
\(\beta(T)\subset\ker\pi=T\); equality follows from
finiteness and the equality of their orders. Thus
\[
                              T\text{ is characteristic in }N.
                                                               \tag{S50.8}
\]
This proof also covers solvable \(L\) or trivial \(Q\).

\paragraph{4. The actual block projectors descend.}
Each \(V_i\) is an irreducible \(T\)-module of dimension
four because its coordinate image is \(L\). Equivalence
of two such modules defines an \(S_r\)-invariant
equivalence relation, so they are either all equivalent or
pairwise inequivalent.

In the first case, choose the common irreducible module \(W\).
The actual evaluation isomorphism is
\[
 W\otimes\operatorname{Hom}_T(W,V_{\mathbf C})
        \longrightarrow V_{\mathbf C},\qquad
                  w\otimes f\longmapsto f(w).
                                                               \tag{S50.9}
\]
Its factor dimensions are four and \(r\). The actual
matrix image of \(T\) is \(\operatorname{End}(W)\otimes I\).
Every element of \(N\) normalizes it. Its action on matrix
units has an inner implementation on \(W\); removing
this implementation leaves an operator commuting with all
matrix units, hence acting only on the multiplicity factor.
Thus (S50.9) is a literal tensor structure for \(N\).
HD-R17 excludes it, since \(4r>16\) and \(N\) is
absolutely irreducible and generated by rank-four elements.

The modules are therefore pairwise inequivalent. Write their
characters as \(\chi_i\) and set \(E=K(\zeta_{\exp T})\).
Character orthogonality gives the literal original projectors
\[
 p_i=\frac4{|T|}\sum_{t\in T}\chi_i(t^{-1})\rho(t),
 \quad p_ip_j=\delta_{ij}p_i,\quad\sum_i p_i=I,
 \quad\operatorname{rank}p_i=4.                            \tag{S50.10}
\]
Their complex images are exactly \(V_i\); they are the
complete primitive central projectors of the given matrix
image of \(E[T]\). Galois automorphisms permute these
projectors and commute with the conjugation action of every
original matrix of \(N\), including its pure permutation
subgroup in the transported coordinates. The centralizer of
this natural \(S_r\) on its labels is trivial. Thus
\[
                  p_i\in M_{4r}(K),\qquad U_i=p_iK^{4r},
        \qquad U_i\otimes_K\mathbf C=V_i.                 \tag{S50.11}
\]
By (S50.8), an actual \(K\)-linear normalizer of \(N\)
normalizes \(T\), so permutes its same primitive central
projectors. This proves the structural assertion.

\paragraph{5. Literal induction, ramified quotient and original corner.}
Return to (S50.1), with the original positive invariant
\(K\)-Hermitian form \(h\), conjugate-linear in the
first argument. The different irreducible \(T\)-blocks
are orthogonal. Equivalently, (S50.10) is self-adjoint under
\(t^\dagger=t^{-1}\), by
\(\chi_i(t^{-1})=\overline{\chi_i(t)}\).
Put \(H=\operatorname{Stab}_G(U_1)\),
\(\sigma=\rho|_{H,U_1}\), and \(H_0=\ker\sigma\).
The restriction \(\sigma\) is absolutely irreducible
because its \(T\)-restriction already is. The actual
induction functions and inverse are
\[
\begin{aligned}
 F_v(x)&=p_1\rho(x)v,&F_v(h_1x)&=\sigma(h_1)F_v(x),\\
 F_{\rho(g)v}(x)&=F_v(xg),&
 F&\longmapsto\sum_{H\backslash G}\rho(x)^{-1}F(x).
\end{aligned}                                                   \tag{S50.12}
\]
The sheet label \(Ht_j^{-1}\leftrightarrow t_jH\)
identifies the right monodromy with the original block action.
These maps give \(V\simeq\operatorname{Ind}_H^G\sigma\).
The induced trace formula then gives
\(K=\mathbf Q(\operatorname{tr}\rho)\subset
\mathbf Q(\operatorname{tr}\sigma)\subset K\), so the
faithful image \(H/H_0\) has exact field \(K\).

From the original connected \(G\)-cover form
\[
 Y=\widetilde C/H,\qquad C'=\widetilde C/H_0,
                 \qquad\alpha:\widetilde C\longrightarrow C'.
                                                               \tag{S50.13}
\]
On an inertia cycle, a fixed section is the successive
transport of a vector fixed by the actual return operator.
Thus (S50.12) identifies the invariant stalks, not only the
unramified local systems. Finite pushforward and this stalk
identification give the actual isomorphism of first parabolic
cohomology with that of the rank-four \(H/H_0\)-system on
\(Y\). Cup products use the sum of the orthogonal block
pairings, which is the finite trace over quotient sheets.

Choose \(v\in U_1(K)\setminus0\), and put \(q=h(v,v)\).
The actual primitive rational idempotents in \(\mathbf Q[G]\)
are equal, because every coefficient off \(H\) vanishes:
\[
\begin{aligned}
 e_v^G&=\frac{2(4r)}{|G|q}\sum_{g\in G}
           \operatorname{Re}h(v,\rho(g^{-1})v)g\\
      &=\frac8{|H|q}\sum_{h_1\in H}
           \operatorname{Re}h(v,\sigma(h_1^{-1})v)h_1=e_v^H.
\end{aligned}                                                   \tag{S50.14}
\]
Averaging over \(H_0\) is the literal embedding of the
same corner from \(\mathbf Q[H/H_0]\). Thus quotient
pullback is an actual isogeny from its primitive factor in
\(JC'\) onto the same integer-cleared image in
\(J\widetilde C\). For covariant left-deck Jacobian
actions, the exact identities are
\[
 \alpha_*\alpha^*=[|H_0|],\qquad
 (\alpha^*)^*\theta_{J\widetilde C}=|H_0|\theta_{JC'},
 \qquad\frac{2|G|}{(4r)q}=\frac{2|H|}{4q}.              \tag{S50.15}
\]
The last equality is the real Schur factor for
\(\operatorname{PD}(j+\bar j)\); it retains the specified
paired Hodge factor and its corner \(K\)-action.

For the originally prescribed line \(Kw\), possibly outside
\(U_1\), use \(A(x)=v h(w,x)/q_w\), \(q_w=h(w,w)\),
and the actual Fourier matrix unit
\[
 b=\frac{2(4r)}{|G|}\sum_{g\in G}
          \operatorname{Re}\operatorname{tr}(A\rho(g^{-1}))g,
                    \qquad b^\dagger b=(q/q_w)e_w.
                                                               \tag{S50.16}
\]
Matrix-coefficient orthogonality gives \(A\), its conjugate,
and zero on the other coefficient factors. An integer multiple
\(Mb\) is an actual \(K\)-linear isogeny between the
primitive images, with pullback polarization multiplier
\(M^2q/q_w\). Together with (S50.15), this applies the
quotient construction to the original specified factor.
A positive rational multiplier changes its rank-six Hermitian
determinant by its sixth power, the norm of its cube.

\paragraph{6. The complete rank-four source endpoint.}
The faithful coefficient system on \(Y\) has rank four,
zero global invariants, and first parabolic dimension six.
Let \(g=g(Y)\), and let \(b\) count exactly its nontrivial
coefficient branch values, with codimensions \(c_j\ge1\).
The exact Euler formula is
\[
                       6=4(2g-2)+\sum_{j=1}^b c_j.
                                                               \tag{S50.17}
\]
For \(g\ge2\) its right side is at least eight, so this
is impossible. For \(g=1\), it gives \(b\le6\), with
pointed-base parameter dimension at most six; here \(b>0\),
so the elliptic-curve automorphism count is the usual pointed
one. For \(g=0\) it gives \(\sum c_j=14\), hence
\(b\le14\). At fixed finite monodromy labels, the covers
over each marked base form a finite discrete set. After a
finite local parameter cover, the actual quotient family
therefore factors through the pointed-curve parameter space.
Values of the original cover that are trivial for this
coefficient system can be forgotten; the intermediate-extension
coefficient variation factors through this same forgetful map.
The actual isogenies (S50.15)--(S50.16), with fixed finite
labels and lattices, preserve the period differential rank.

Rank nine consequently forces \(g=0\) and
\[
                       12\le b\le14,\qquad\sum c_j=14.
                                                               \tag{S50.18}
\]
There are at least ten reflection inertia values, but there
can also be nonreflection inertia: (S50.18) is not replaced
by an all-reflection statement. Apply HD-S40's complete
rank-at-most-fourteen source theorem to this actual faithful
rank-four quotient group, with exact field \(K\), prescribed
primitive corner, restricted theta polarization, signature
\((3,3)\), and period rank nine. It gives precisely the
already supplied determinant classes stated above. The
positive rational multipliers in (S50.15)--(S50.16) preserve
these ordinary determinant classes. Compatible flat changes
of polarization have the exact scope of HD-S17.

The proof closes complex block dimension four for an
absolutely irreducible \(N_4\) in the remaining \(m\ge21\)
range, since its divisibility by four gives \(r\ge6\).
It does not claim descent of the separate block dimensions
two and three. No finite classification or numerical experiment
is substituted for the actual kernel and quotient arguments.
